\chapter{验证、不确定性与可辩护的决策}\label{ch:practice}

\begin{chapterguide}
\guidequestion{一个磁轭模型的网格已经加密过，线性方程组的残差小到 $10^{-13}$，场图也平滑好看。凭它就能决定磁轭板做多厚、梯度线圈的电流上限定多少吗？如果不能，还缺哪几种证据？}
\begin{guidegoals}
\item 从要回答的物理问题和要做的决定出发，选择观测量、近似、表述与检查，而不是从软件模块出发；
\item 区分数值验证与物理确认，说出每一种证据能发现什么错误、发现不了什么；
\item 按应用列出必须通过的检查，认出全书反复出现的几类失误；
\item 为一个观测量做误差预算：一次只改变一个数值环节，独立扰动不确定的物理输入，判断哪一项主导；
\item 写出一份别人能复现、也能追问的结果记录，并说清教学模型在哪里停止。
\end{guidegoals}
\guideprereq{\secref{sec:femscalar-convergence}（收敛实验，残差与离散误差）；\secref{sec:lumped-validity}（网络与场解的对照）；\chref{ch:framework}（三族算子与观测量）；\chref{ch:projects}（模型约定与综合项目）。}
\end{chapterguide}

\section{从物理问题选择流程}\label{sec:practice-workflow}

先回到\chref{ch:magnets} 的二维磁体截面。网络给出的气隙中心磁场是 $0.191\,\mathrm T$，场解给出 $0.152\,\mathrm T$；场解的网格加密一倍，中心磁场变化不到 $0.1\%$，外边界从 $1.5\,\mathrm m$ 扩大到 $6\,\mathrm m$，它也几乎不变。这些检查说明了什么？它们说明：\emph{这个二维模型的方程解得很准}。它们不能说明一台三维磁体的中心磁场是 $0.152\,\mathrm T$，也不能说明真实材料的磁轭会不会饱和——那取决于三维几何、实测的 B--H 数据和装配公差。“算得准”与“算得对”是两件事；而设计决定需要的，是知道“对到什么程度”。

本章把这几件事分开，并说明怎样为一个设计决定收集足够的证据。第一步，是让流程服从问题。

\subsection{六个问题}
面对一个新的电磁仿真任务，按下面的顺序回答六个问题。顺序很重要：每一个问题的答案都限制了后面的选择。
\begin{enumerate}
\item \textbf{要做什么决定？}由此确定观测量。观测量不是“算出场”，而是“目标区域的平均电场”“成像区的峰峰值均匀度”“磁轭中最大的 $B$”“外周神经的兴奋阈值”这样具体的读出（\figref{fig:field_observables}）。
\item \textbf{哪些物理必须保留？}按最快的上升沿和相应的长度，用三把尺子 \eqref{eq:scales} 检查。组织在 $3\,\mathrm{kHz}$ 下 $\eta_m\approx4.7\times10^{-4}$，可以单向耦合；同一个人体在 $0.1\,\mathrm T$ 的射频频率下 $\eta_m\approx0.67$，就不能先验地这样做（\chref{ch:rf}）。
\item \textbf{真正知道哪些本构信息？}组织电导率的范围、B--H 数据的来源与适用范围、剩磁随温度的变化。不知道的东西要写成不确定性，而不是一个看起来很确定的数。
\item \textbf{源怎样给定、怎样归一化？}电极电流、线圈电流的变化率、剩磁、端口功率——单位场图的含义取决于它（\chref{ch:tms}）。
\item \textbf{用什么表述与未知量？}三族算子中的哪一族（\figref{fig:application_map}），用 $\phi$、$\vb A$ 还是 $\psi$，线性还是非线性。
\item \textbf{什么证据能让结论站得住？}这是本章其余各节的内容。
\end{enumerate}

\begin{intuition}[先问“要决定什么”，再问“要算什么”]
同一个磁体，问“气隙平均磁场多大”时，网络模型几毫秒就能给出可用的估计；问“磁体尖角会不会退磁”时，再精细的网络也无能为力，必须做非线性场计算，并且要有工作温度下的退磁曲线。求解器、网格策略和源的表示都应当服从问题，而不是反过来。
\end{intuition}

\begin{example}[label=ex:practice-workflow]{同一台 0.1 T 永磁 MRI 上的三个问题}
一台 $0.1\,\mathrm T$ 的永磁 MRI。(a) 磁轭会不会局部饱和？(b) 梯度线圈切换时，病人体内的电场会不会引起外周神经刺激？(c) 射频线圈在病人体内沉积多少功率？分别说出观测量、需要保留的物理和至少一项关键检查。
\solution
(a) 观测量是磁轭中局部的 $|\vb B|$ 与磁体内的反向场，不是平均气隙磁场。静磁问题，要用非线性 B--H 曲线，并用一致切线的牛顿迭代求解（\chref{ch:femvector}）；关键检查：网格与外边界的收敛、$|\vb B|$ 是否超出 B--H 数据的范围、对称面条件是否与源的奇偶一致（\chref{ch:magnets}）。

(b) 观测量是沿神经路径的驱动与阈值，不是最大 $|\vb E|$（\chref{ch:neuro}）。组织中 $\eta_m$ 很小，可用单向耦合的指定源感应（\chref{ch:tms}）；但极靴、磁轭若导电，它们中的涡流会改变梯度场本身，那一部分要用耦合 MQS（\chref{ch:eddy}），而且是在被磁体偏置的磁轭上线性化 \eqref{eq:biaslinearization}。关键检查：单位场图的归一化（每 A/s 还是每 A/$\mu$s）。

(c) 观测量是在规定质量与时间上平均的 SAR。$4.26\,\mathrm{MHz}$ 下 $\eta_m\approx0.67$，不能直接搬用梯度线圈的单向耦合模型，要用该频率下的材料数据并与更完整的模型对照（\chref{ch:rf}）；关键检查：比较 SAR 之前先把源归一化到同一个激发目标（\secref{sec:rf-normalization}）。

三个问题来自同一台机器，却是三个不同的电磁问题。
\end{example}

\section{数值验证与物理确认}\label{sec:practice-vv}

\subsection{两个问题}
有两个容易混在一起的问题\cite{oberkampf}：
\begin{itemize}
\item \textbf{数值验证}（verification）：离散模型是否正确地解了\emph{写下的}方程？
\item \textbf{物理确认}（validation）：写下的方程和输入，是否足以代表真实的装置或实验？
\end{itemize}
网格收敛、对称性、残差、与解析解的比较，都属于前者；与场探头、阻抗测量、MRI 场图、刺激阈值实验的比较，属于后者。视觉上合理的场图，两样都不自动满足。

\fig{verification_validation}{从真实系统到数值结果的建模链。中间一列自上而下：真实系统经过理想化成为概念模型（决定保留哪些物理），写成数学模型（方程、边界条件、源与本构），离散成矩阵方程，编成程序，求解得到数值结果。右侧：数值验证只检查下面几步——方程有没有解对；上面两步的建模假设它看不见。左侧：物理确认把数值结果与真实系统的测量放在一起比较，比较的必须是同一个观测量。}{fig:verification_validation}

\figref{fig:verification_validation} 画出了两者作用的位置。从真实系统到数值结果，要经过五步：理想化、写成方程、离散化、编程、求解。数值验证检查的是后三步：离散方程是否逼近了数学模型，程序是否实现了离散方程，求解是否足够准。前两步——保留了哪些物理、略去了哪些耦合、用了什么材料数据——数值验证完全看不见：一个解得极准的错误方程，照样能通过所有网格收敛检查。只有把结果与真实系统的测量放在一起，才能检查这两步。

\begin{compare}{数值验证与物理确认}{数值验证（verification）}{物理确认（validation）}
\cmprow{问的问题}{离散模型是否正确地解了写下的方程}{写下的方程和输入是否足以代表真实装置}
\cmprow{比较的对象}{数值解与精确解、加密解、守恒律}{数值结果与测量（同一个观测量）}
\cmprow{典型证据}{制造解的收敛阶；网格序列；残差；对称与守恒}{场探头、阻抗、MRI 场图、刺激阈值实验}
\cmprow{能发现}{实现错误、网格不足、边界条件施加错误}{缺少的物理、错误的材料数据、源标定错误}
\cmprow{发现不了}{方程本身选错了}{数值误差与模型误差碰巧相互抵消}
\cmprow{次序}{先做：未经验证的解拿去确认，没有意义}{后做：只在验证过的范围内有意义}
\cmpsum{验证问“方程解对了吗”，确认问“方程选对了吗”。一张好看的场图，两样都不自动满足。}
\end{compare}

\subsection{验证本身也分两种}
\textbf{代码验证}检查程序：它实现的算法对不对。最干净的办法是制造解（\secref{sec:femscalar-convergence}）：先选定一个解，算出它需要的源，再看数值解是否以理论阶数收敛。线性三角形单元实测的阶数是 $1.998$（电势）与 $0.999$（梯度）。\textbf{解的验证}检查这一次计算：对\emph{这个}问题、\emph{这个}观测量，离散误差有多大。它需要一串加密的网格。\chref{ch:magnets} 的磁体截面上，网格加密一倍，中心磁场变化不到 $0.1\%$，这是一项通过了的解的验证；\chref{ch:tes} 中等势电极边缘的峰值在三张网格上依次是 $6.9$、$9.8$、$14$，不断增大——这项验证\emph{失败}了，它说明“理想电极边缘的点峰值”不是一个定义良好的观测量（\secref{sec:tes-electrodes}）。

\begin{compare}{代码验证与解的验证}{代码验证}{解的验证}
\cmprow{对象}{程序：它实现的算法对不对}{这一个具体问题的解：误差有多大}
\cmprow{需要}{一个已知的精确解（常用制造解）}{一串逐次加密的网格，或可靠的误差估计}
\cmprow{判据}{实测收敛阶等于理论阶}{所关心的观测量在加密下稳定}
\cmprow{本书的例子}{\secref{sec:femscalar-convergence}：阶数 $1.998$ 与 $0.999$}{\chref{ch:magnets}：网格加密一倍，中心磁场变化小于 $0.1\%$}
\cmprow{不能替代}{解的验证：程序对，网格仍可能太粗}{代码验证：同一个程序错误在每张网格上都在}
\cmpsum{代码验证一次性地证明“程序会解方程”；解的验证每一次都要证明“这一次解得够准”。}
\end{compare}

\subsection{介于两者之间：独立的表述与模型之间的对照}
另有两类常用的检查，既不是纯粹的验证，也不是确认。
\begin{itemize}
\item \textbf{独立的表述。}同一个量用两种数学上等价、实现上独立的方法算。\chref{ch:framework} 中饱和铁芯的吸力，由共能求导、由能量求导、由 Maxwell 应力，三种算法都给出 $405\,\mathrm N$；把割线电感代入线性公式只得到 $31.2\,\mathrm N$。这类检查能抓住实现与推导中的错误，但三种算法共用同一组建模假设，对确认没有帮助。
\item \textbf{模型之间的对照。}\secref{sec:lumped-validity} 用二维场解检查磁路网络：气隙中心的 $B$ 相差不到 $1.5\%$，绕组铁柱中的磁通却在 $1\,\mathrm{mm}$ 气隙时差 $26\%$、$8\,\mathrm{mm}$ 时差到 $2.50$ 倍。这检验的是网络的\emph{压缩假设}，参照物是一个更完整、并且已经验证过的模型；它不是对真实装置的确认。
\end{itemize}

\subsection{记录本身也是证据}
别人要能检查你的验证，就要知道你算的是什么。一份可靠的记录至少写明：几何的来源、网格统计、单元类型、本构关系、边界条件、求解器容差、源的归一化，以及图中画的到底是哪个物理量。非线性磁场还要说明：报告的是割线还是微分磁导率，用了哪一条 B--H 分支，磁滞是忽略了还是怎样近似的（\chref{ch:bh}）。完整的清单在\secref{sec:practice-reproducible}。

\section{各应用的检查清单}\label{sec:practice-checklists}

前面各章在讲每个应用时，都留下了一些“必须检查”的事。\tabref{tab:practice-check-fields} 与\tabref{tab:practice-check-responses} 把它们收在一起，每一项都注明了讲解它的地方。它们不是额外的规则，而是各章推导的直接推论；拿不准某一项为什么要查时，回到出处去看。

\begin{table}[htbp]\centering\small
\caption{场问题的检查清单。}\label{tab:practice-check-fields}
\begin{tabularx}{\linewidth}{@{}>{\raggedright\arraybackslash}p{.13\linewidth} Y@{}}
\toprule
应用 & 必须检查什么（出处）\\\midrule
TES & ① 电极电流之和为零，求解后在\emph{所有}节点上检查 $\vb K\boldsymbol\phi-\vb f$（\eqref{eq:currentcompatibility}，\exref{ex:tes-ground}）；② 固定参考不改变场，把整个表面接地是另一个实验（\secref{sec:tes-boundary}）；③ 电导率张量的坐标约定 \eqref{eq:sigmatensor}；④ 电极附近的点峰值是否随网格增长，必要时改用完整电极模型或区域平均（\secref{sec:tes-electrodes}）；⑤ 电源固定的是电流还是电压\\
TMS & ① 源线性：电流变化率加倍，电场加倍，两个线圈的场相加（\secref{sec:tms-operator}）；② 单位场图是每 A/s 还是每 A/$\mu$s，匝数只计入一次 \eqref{eq:tmswaveform}；③ 绝缘边界是 $\vb n\cdot\sig(\vb E_p-\nabla\phi)=0$，不是 $\partial_n\phi=0$，并检查正交性残差 \eqref{eq:tmsorthogonal}；④ 初级场不按组织标签裁剪，解析源与有限元右端用同一个规范（\secref{sec:potentials-gauge}）；⑤ 线圈源单独与圆环、螺线管的解析式对照（\eqref{eq:loopaxis}、\eqref{eq:solenoidA}）\\
MRI 梯度 & ① 设计区域内的空间线性，以及 Maxwell 方程带来的伴随场（\eqref{eq:roiMaxwell}、\eqref{eq:concomitant}）；② 电感、端电压与储能 \eqref{eq:circuit}；③ 绕组在主磁场中受的力；④ 附近导体中的涡流与时间响应（\chref{ch:eddy}）；⑤ 体内的感应电场与外周神经刺激（\chref{ch:neuro}）\\
匀场 & ① 区分恒定偏移、低阶谐波与随时间的漂移（\chref{ch:mri}）；② 均匀度指标的定义：RMS 还是峰峰值，在什么区域、按什么权重\\
磁体与磁轭 & ① 磁体内的反向场与工作温度下的拐点（\secref{sec:bh-hysteresis}）；② 磁轭的局部饱和，不只看平均值（\tabref{tab:magnets-outputs}）；③ 空气域与外边界：扩大后比较所需的观测量，杂散场需要大得多的区域（\chref{ch:magnets}）；④ 对称面条件按源的奇偶选择（\figref{fig:symmetry_bc}）；⑤ 力在空气中的闭合面上积分，并改变积分面位置检查 \eqref{eq:maxwellstress}；⑥ 集总模型先标定、后优化（\secref{sec:lumped-validity}）\\
\bottomrule
\end{tabularx}
\end{table}

\begin{table}[htbp]\centering\small
\caption{响应与读出的检查清单。}\label{tab:practice-check-responses}
\begin{tabularx}{\linewidth}{@{}>{\raggedright\arraybackslash}p{.13\linewidth} Y@{}}
\toprule
读出 & 必须检查什么（出处）\\\midrule
神经响应 & ① 场投影到明确的神经路径，用路径积分构造节点驱动 \eqref{eq:nervedifference}；② 场插值、路径采样、电缆离散与时间步分别加密（\chref{ch:projects}）；③ 兴奋判据先写清，二分搜索前先扫描，确认响应在区间内单调（\chref{ch:neuro}）；④ 用实际的波形；强度--时程曲线不在设备之间直接搬用\\
加热（SAR） & ① 峰值还是有效值（\secref{sec:framework-sar}）；② 质量平均的区域、空气的处理与时间平均（\secref{sec:rf-averages}）；③ 多通道保留交叉项 $\vb c^\dagger\vb Q\vb c$ \eqref{eq:sarQ}；④ SAR 不是温度，也不是神经响应\\
射频 & ① 按载频重新量一遍三把尺子 \eqref{eq:scales}；② 端口归一化，比较之前先对齐激发目标（\secref{sec:rf-normalization}）；③ 外边界用吸收条件或真实存在的屏蔽，不能随手令 $\widetilde{\vb E}=0$；④ 旋转分量的时间约定与旋向 \eqref{eq:corotating}\\
\bottomrule
\end{tabularx}
\end{table}

\section{常见失误}\label{sec:practice-failures}

仿真可以内部完全一致，却在物理上误导人。回头看全书，最常见的失误多半不是求解器算错了，而是约定、归一化或观测量出了问题：程序忠实地解了一个与本意不同的问题。\tabref{tab:practice-failures} 把它们按“症状—怎样查”整理在一起；每一行的例子都在前面某一章里真正算过。

\begingroup\small\renewcommand{\arraystretch}{1.22}
\begin{longtable}{@{}>{\raggedright\arraybackslash}p{.25\linewidth} >{\raggedright\arraybackslash}p{.39\linewidth} >{\raggedright\arraybackslash}p{.3\linewidth}@{}}
\caption{全书出现过的失误。}\label{tab:practice-failures}\\
\toprule
失误 & 全书中的例子 & 怎样查\\\midrule\endfirsthead
\toprule
失误 & 全书中的例子 & 怎样查\\\midrule\endhead
峰值与有效值混用 & 同一个“$100\,\mathrm{V/m}$”给出 $2.5$ 或 $5\,\mathrm{W/kg}$（\secref{sec:framework-sar}） & 读入任何场数据之前先问它的约定\\
单位场图差 $10^6$，或匝数计入两次 & 把每 A/$\mu$s 的场图当作每 A/s 用（\chref{ch:tms}） & 用一个能手算的算例（长螺线管）核对量级\\
只比较依赖规范的量 & 两种规范的 $-\partial_t\vb A$ 截然不同，总电场却相同（\exref{ex:potentials-gauge}） & 只比较 $\vb E$、$\vb B$ 这样的规范不变量\\
把 $|\vb E|$ 当作神经兴奋的充分指标 & 局域场中 $|\vb E|$ 最大处激活函数恰为零（\figref{fig:activating_function}） & 投影到神经路径，解电缆方程\\
把局部斜率当作整个部件的磁导率 & 冻结割线把偏置附近的小信号响应高估 $12.1$ 倍（\chref{ch:femvector}） & 用一致切线，并用一串步长的中心差分检验 \eqref{eq:biascheck}\\
非线性磁体的空气域太小 & $1\,\mathrm m$ 处的杂散场随外边界从 $0.97$ 变到 $0.68\,\mathrm{mT}$（\chref{ch:magnets}） & 扩大外边界，比较\emph{所需的}观测量\\
参考节点吞掉不平衡的电流 & 写错的 $-0.9\,\mathrm{mA}$ 被悄悄“改”成 $-1\,\mathrm{mA}$（\exref{ex:femscalar-twolayer}） & 求解前查 $\vb 1^{\mathsf T}\vb f$，求解后查所有行的残差\\
对称面条件只看外形 & 磁体的四分之一模型不能原样用来算梯度场（\figref{fig:symmetry_bc}） & 按源在镜像下的奇偶选条件\\
把单通道 SAR 图相加 & 同相发射时中点 SAR 是两张图之和的两倍（\figref{fig:sar_cross_terms}） & 先合成复数场，或保存 $\vb Q$ 矩阵\\
把随网格增长的点峰值当作热点 & 等势电极边缘的峰值 $6.9$、$9.8$、$14$（\secref{sec:tes-electrodes}） & 圆角、接触层，或在规定体积上平均\\
用方便的数学惩罚代替真实约束 & 系数的平方和不是功率 $\vb c^{\mathsf T}\vb R\vb c$（\chref{ch:coils}） & 目标函数的每一项都要对应一个物理量\\
未经验证的模型直接进优化器 & 在 $1\,\mathrm{mm}$ 气隙处标定的网络，到 $8\,\mathrm{mm}$ 时把柱磁通低估 $48\%$（\secref{sec:lumped-validity}） & 先在设计空间上标定，把设计变量限制在验证过的范围内\\
\bottomrule
\end{longtable}
\endgroup

表中的失误有一个共同点：它们都不会让程序报错。求解器只负责解它拿到的方程；约定、归一化与观测量是不是本意，只有建模者自己能检查。

\section{误差预算}\label{sec:practice-errorbudget}

\subsection{误差从哪里来}
一个数值结果与真实值之间的差，来自许多互不相同的环节（\tabref{tab:practice-sources}）。降低线性方程组的残差，只减小其中一项；加密网格，也只减小其中一两项。

\begin{table}[htbp]\centering\small
\caption{误差与不确定性的来源。}\label{tab:practice-sources}
\begin{tabularx}{\linewidth}{@{}>{\raggedright\arraybackslash}p{.13\linewidth} Y Y >{\raggedright\arraybackslash}p{.13\linewidth}@{}}
\toprule
来源 & 本书中的例子 & 怎样估计 & 加密网格能减小吗\\\midrule
几何与分割 & 网格化时消失的薄层（\exref{ex:femscalar-thinlayer}）；按单元标签近似的夹杂 & 扰动界面位置；与原始影像或图纸对照 & 不能\\
材料 & 组织电导率、B--H 曲线、剩磁及其温度依赖 & 在有依据的范围内独立扰动；记录数据来源 & 不能\\
源与标定 & 线圈波形、电极接触、端口归一化 & 测量；与解析的源场对照 \eqref{eq:loopaxis} & 不能\\
模型约化 & 单向耦合、准静态、集总化 & 三把尺子；与保留更多耦合的模型对照 & 不能\\
离散 & 网格尺寸、单元阶次、几何的表示 & 网格序列；制造解 & 能\\
代数求解 & 迭代容差、非线性迭代的步数 & 收紧容差，看残差与观测量 & 收紧容差能\\
插值与后处理 & 跨界面平均矢量、对噪声场二次求导、体素采样 & 换一种插值；沿路径积分 & 部分\\
响应模型 & 膜模型、兴奋判据、热模型 & 与实验对照；换模型比较 & 不能\\
\bottomrule
\end{tabularx}
\end{table}

\subsection{一次只改一处}
做误差预算有两条基本纪律：
\begin{itemize}
\item \textbf{数值环节一次只改一个。}保持物理问题完全不变——同样的几何、材料、源——只加密网格，或只收紧容差，或只换插值方式。同时改两样，就分不清变化来自哪里。
\item \textbf{不确定的物理输入独立扰动。}在有依据的范围内，一次只改变一个输入，看观测量变化多少。由此才知道哪一项主导。
\end{itemize}
报告的应当是\emph{观测量} $y$ 的变化，而不只是全局的场范数。不同的观测量收敛得快慢不同：局部峰值、体积平均、磁链、神经阈值，可以有完全不同的收敛速度。理想尖角处的点峰值甚至可能随加密一直增长；这时应当改用物理上规定的有限区域、圆角几何或其他定义良好的观测量，而不是把那个峰值当作一个与网格无关的材料性质。

对远离零的正量，相邻两张网格之间的相对变化是一个方便的指标：
\begin{keyeq}
\epsilon_h=\frac{|y_{h/2}-y_h|}{|y_{h/2}|}\qquad\text{（$y$ 远离零）};\qquad
\epsilon_h=\frac{|y_{h/2}-y_h|}{y_{\rm ref}}\qquad\text{（$y$ 接近零）}.
\label{eq:practice-indicator}
\end{keyeq}
$y$ 接近零时（例如对称面上本应为零的一个场分量），相对指标会因为分母很小而失去意义，要换成一个明确定义的绝对参考尺度 $y_{\rm ref}$，例如该区域中场的典型大小。

\subsection{一个算例}

\begin{example}[label=ex:practice-budget]{同一个观测量，两种刺激的误差预算}
取\figref{fig:tes_setup} 的含夹杂导电方块，把夹杂当作“目标组织”，观测量 $y$ 是夹杂内 $|\vb E|$ 的面积平均。对 TES（恒流、均匀电流密度电极）与 TMS（均匀的 $\partial_tB_z$）分别做下面的改变，一次只改一项，以每边 $n=128$ 段的网格为基准：网格 $n=64$、$128$、$256$；把直接法换成共轭梯度法，相对残差容差 $10^{-3}$；夹杂半径 $0.16\to0.15$、$0.17$；夹杂电导率 $10\to8$、$12$；（TES）电极模型。扰动的幅度都是教学输入。哪一项主导？
\solution
\emph{数值误差。}TES 的 $y$ 在三张网格上为 $0.1335$、$0.1301$、$0.1291$，相邻网格的变化 \eqref{eq:practice-indicator} 为 $2.6\%$ 与 $0.80\%$；TMS 为 $0.0536$、$0.0540$、$0.0540$，变化 $0.62\%$ 与 $0.12\%$。直接法的相对残差为 $1.7\times10^{-13}$；容差放宽到 $10^{-3}$，TES 的 $y$ 也只变 $0.02\%$，TMS 只变 $0.002\%$。

\emph{模型输入。}TES：夹杂电导率 $\pm20\%$ 使 $y$ 变化 $+22.5\%$、$-15.5\%$；夹杂半径变化 $\pm0.01$，$y$ 变化 $+1.4\%$、$-0.4\%$；把均匀电流密度电极换成等势电极，$y$ 降低 $8.5\%$，换成完整电极模型（$z\sigma/w=0.08$、$0.8$）降低 $5.2\%$、$1.5\%$。TMS：同样 $\pm20\%$ 的电导率只改变 $+0.78\%$、$-0.45\%$；半径减小 $0.01$ 时 $y$ 却降低 $6.1\%$，增大 $0.01$ 时升高 $6.1\%$。

\emph{结论。}两种刺激在 $n=128$ 时的离散误差都已小于 $1\%$，求解器容差完全不重要。TES 的这个观测量由夹杂电导率主导：夹杂内的电场取决于流进它的电流与它自己的电导率，高电导率的夹杂里电场被压低，电导率稍有不同，电场就跟着明显改变。TMS 的同一个观测量几乎不在乎夹杂电导率。原因要看夹杂里的场由什么组成：均匀 $\partial_tB_z$ 的初级场，相对夹杂自己的中心可以拆成两部分。一部分绕夹杂中心转圈，处处与夹杂边界相切，不在边界上留下电荷，所以与夹杂电导率无关；另一部分是一个均匀场，它要穿过夹杂边界，被高电导率的夹杂压到原来的不足五分之一，只有这一部分随电导率变化。它本来就小，又与转圈的部分在各处方向不同，$\pm20\%$ 的变化只让平均 $|\vb E|$ 改变不到 $1\%$。同一个观测量却对分割很敏感：转圈的那部分离夹杂中心越远越强，平均值正比于半径；这里的观测量又是在“夹杂这块区域”上平均的，半径一变，平均的区域和其中的场就一起变了。
\end{example}

\fig{error_budget}{\exref{ex:practice-budget} 的误差预算：夹杂内平均 $|\vb E|$ 的相对变化（对数横轴，虚线为 $1\%$）。蓝灰色是数值误差（网格、求解器容差），沙色是模型输入的不确定性（几何、材料、电极模型）。同一组扰动，TES 由材料主导，TMS 由分割主导；两者的离散误差在 $n=128$ 时都小于 $1\%$。扰动幅度为教学输入。}{fig:error_budget}

这个算例说明了误差预算最重要的一点：\emph{它属于“一个问题加一个观测量”，不属于软件或网格。}同样 $\pm20\%$ 的电导率不确定性，对 TES 的这个观测量是两成，对 TMS 的同一个观测量不到百分之一。在 TES 中，网格从 $128$ 加密到 $256$ 只改变 $0.80\%$，而材料带来 $\pm20\%$ 量级的不确定——再加密网格，结论不会变；该花力气的是电导率的数据。

还要注意，“全部电导率乘以同一个因子”这种不确定性有精确的答案：恒流 TES 中电场与该因子成反比，TMS 中电场完全不变（\chref{ch:tes}、\chref{ch:tms}）。检查脚本对本例确认了这两条。

\begin{compare}{数值误差与模型不确定性}{数值误差}{模型不确定性}
\cmprow{来源}{离散、代数求解、插值：方程没有解准}{几何、材料、源、近似、响应模型：方程或输入不准}
\cmprow{怎样估计}{网格序列、制造解、收紧容差、守恒检查}{独立扰动输入；与更完整的模型或测量对照}
\cmprow{能否靠计算消除}{原则上能：加密、提高阶次、收紧容差}{不能：算得再多也不知道真实的电导率}
\cmprow{\exref{ex:practice-budget}（TES）}{网格 $128\to256$：$0.80\%$；容差 $10^{-3}$：$0.02\%$}{夹杂电导率 $\pm20\%$：$+22.5\%$、$-15.5\%$}
\cmprow{\exref{ex:practice-budget}（TMS）}{$0.12\%$；$0.002\%$}{夹杂半径 $\pm0.01$：$\pm6.1\%$}
\cmprow{报告方式}{给出加密序列和实测的收敛}{写明扰动的范围及其依据}
\cmpsum{数值误差回答“方程解得准不准”，可以算到足够小；模型不确定性回答“方程与输入对不对”，只能估计与缩小。前者小于后者以后，再加密网格就不会改变结论。}
\end{compare}

怎样把各项合成一个总的不确定度，需要明确的统计假设（各项是否独立、服从什么分布），超出了本书的范围。即使不做合成，“哪一项主导、它有多大”这两个答案，往往已经足以支持或否定一个设计决定。

\subsection{两张网格一致，不等于收敛}
相邻网格的 $\epsilon_h$ 很小是证据，但它不是真实误差的上界。

\begin{example}[label=ex:practice-twomesh]{两张网格一致，两张都错}
\exref{ex:twolayer} 的双层导体（总厚 $10\,\mathrm{mm}$，其中 $2\,\mathrm{mm}$ 的低电导层，真实电阻 $1400\,\Omega$），用 1 到 10 个等长单元离散，每个单元按其中心所在的层取电导率（\exref{ex:femscalar-thinlayer} 的做法）。1 个与 2 个单元的结果一致吗？能据此说网格已经收敛吗？
\solution
1 个与 2 个单元都给出 $500\,\Omega$，$\epsilon_h=0$：两张网格完全一致。可是两者都比真实值低 $64\%$——所有单元的中心都落在第 1 层，低电导层在两张网格上都消失了。继续加密，3 到 7 个单元依次给出 $2000$、$1625$、$1400$、$1250$、$1143\,\Omega$，忽高忽低；只有界面恰好落在节点上（5 个、10 个单元）时才精确。每一次计算的代数残差都小于 $10^{-12}$。网格必须先\emph{表达}几何，收敛检查才有意义；加密时还要让几何的表示一起改进，而不是只把同一个错误的几何切得更细。
\end{example}

所以，制造解、独立的表述、守恒检查、网格序列与实验标定，各自约束的是不同的失误。任何一种单独都不够，好在它们也很少同时失败。

\begin{pitfall}[“残差这么小，结果一定很准”]
代数残差只说明“按这张网格写下的方程”解得准。\exref{ex:practice-twomesh} 中残差小于 $10^{-12}$，电阻却错了 $64\%$；\secref{sec:femscalar-convergence} 的制造解中残差约 $10^{-13}$，电场的相对误差仍有 $2.45\%$。残差、离散误差与模型不确定性是三个不同的量，要分别给出证据。
\end{pitfall}

\subsection{误差预算与决定}
本章标题里的“可辩护”，指的就是这一步：把观测量的不确定范围与做决定所需的裕量放在一起比较。若磁轭的局部 $B$ 在所有可信的材料曲线与网格下都远离拐点，就可以放心地减薄；若它恰好落在拐点附近，结论就取决于 B--H 数据本身（\secref{sec:lumped-validity} 末尾的误区），这时该做的是补充材料数据或测量，而不是再加密网格。一个设计决定能被辩护，不是因为某个数“算得很精确”，而是因为说得出：哪些误差已经小到无关紧要，哪一项主导，它为什么不会改变结论。

\section{可复现的结果}\label{sec:practice-reproducible}

\subsection{随结果一起保存什么}
一个结果能不能被别人复现、检查和追问，取决于它带着哪些信息。\tabref{tab:practice-record} 列出了最低限度的内容。其中第一行，就是\chref{ch:projects} 中贯穿所有项目的模型约定（\secref{sec:projects-contract}）。

\begin{table}[H]\centering\small
\caption{一个可复现的结果随数据一起保存的内容。}\label{tab:practice-record}
\begin{tabularx}{\linewidth}{@{}>{\raggedright\arraybackslash}p{.2\linewidth} Y Y@{}}
\toprule
内容 & 包括 & 缺了它会怎样\\\midrule
模型约定 & 几何与坐标单位、源及归一化、材料定律、保留的 Maxwell 项、未知量、边界与界面、参考或规范、观测量与平均方式（\secref{sec:projects-contract}） & 别人解的是另一个问题\\
来源 & 几何、分割与材料数据的出处和版本；B--H 曲线与剩磁的来源及温度 & 无法判断输入是否可信\\
源与波形 & 确切的电流或电压波形、端口归一化、匝数 & 单位场图无法换算成真实的场\\
坐标系 & 网格、线圈位姿、神经路径与材料轴所在的公共坐标系 & 旋转与平移出错却看不出来\\
数值设置 & 网格统计、单元类型与阶次、求解器与容差、非线性迭代的判据 & 无法重做验证\\
收敛与检查记录 & 网格序列上的观测量、守恒与解析对照的结果 & 别人只能“相信”\\
后处理脚本 & 把场变成报告指标的程序（插值、投影、平均） & 同一张场图读出不同的数\\
偏置状态 & 非线性问题的工作点与对应的切线（导出小信号场图时） & 小信号场图离开了它的有效范围\\
\bottomrule
\end{tabularx}
\end{table}

两条容易被忽视的要求：
\begin{itemize}
\item \textbf{单位与归一化要跟着数组走}，写在数据文件的元数据里，而不是只写在笔记的某一段。场数组不会自己解释自己（\secref{sec:projects-interface}）。
\item \textbf{导出小信号场图时，同时保存非线性偏置状态和对应的切线} \eqref{eq:biaslinearization}。它们只在那个偏置附近有效；离开了偏置的说明，这张图就无法正确使用。
\end{itemize}

\subsection{本书自己的做法}
本书正文中的每一个数字，要么来自上一版讲义，要么由配套的检查脚本 \texttt{checks/check\_chNN.py} 计算，并用断言检验；插图模块与检查脚本导入的是同一套模型代码。所以正文、插图与断言共用同一份实现：改动一个模型而忘了改正文，脚本就会报错。这就是\tabref{tab:practice-record} 中“后处理脚本”与“检查记录”两行在小尺度上的样子。

这些脚本检查的是：标量有限元的收敛基准、解析的概念算例、导体的矩阵、源的归一化、射频的旋转分量与 SAR 的代数。它们\emph{不是}经过验证的三维人体模型，不是全波射频求解器，也不是临床阈值的预测器。把\chref{ch:projects} 的项目用到真实的装置上，还需要真实的几何与材料，以及对所用软件的验证。知道一个教学模型到哪里为止，本身就是建立可信模型的能力。

\section{下一步学什么}\label{sec:practice-next}

\subsection{一条迭代的路线}
学习不必一步走到三维头模型或整台磁体。更可靠的路线是迭代的：先做能解析检查的一维、二维问题，用有限元复现它们；再一点点加入真实的几何、材料与源。目标不只是学会操作软件，而是形成判断：结果中的哪些特征由 Maxwell 方程决定，哪些来自材料假设，哪些只是数值近似。

\subsection{本书到哪里为止}
本书的综合项目为建立和检验真实模型打下基础，但不能代替经过验证的解剖数据、经过标定的线圈波形、实测的 B--H 数据、生产级的三维或全波求解器，以及临床与标准层面的确认。高级的磁滞模型、热生理的耦合、网格生成算法、完整的射频网络设计与神经元群体建模，都是更进一步的专题；把它们压缩成缺少适用条件的操作配方，只会误导。\tabref{tab:practice-next} 给出几个方向的下一步。

\begin{table}[htbp]\centering\small
\caption{本书到哪里为止，下一步需要什么。}\label{tab:practice-next}
\begin{tabularx}{\linewidth}{@{}>{\raggedright\arraybackslash}p{.14\linewidth} Y Y >{\raggedright\arraybackslash}p{.2\linewidth}@{}}
\toprule
方向 & 本书给出了什么 & 还需要什么 & 可以从哪里开始\\\midrule
TES 与 TMS 头模型 & 电导方程、完整电极模型、导联场、单位场图、投影检查 & 经过验证的分割与电导率数据，头部网格的生成 & SimNIBS~\cite{simnibs}；组织参数~\cite{itislf,itisdielectric}\\
磁体与磁轭 & 磁路、二维 $A_z$、增量磁阻率张量、牛顿迭代、应力张量 & 实测的 B--H 与退磁曲线及其温度依赖，三维棱单元求解器 & FEMM~\cite{femm}（二维）；MFEM~\cite{mfem}、GetDP~\cite{getdp}\\
线圈设计 & 流函数、场、电阻与电感算子、正则化 & 真实绕线与制造约束、热设计、与测量对照 & bfieldtools~\cite{bfieldtheory,bfieldsoftware}\\
射频 & 频域方程与弱形式、旋转分量、SAR 二次型与三种平均 & 生产级全波求解器，端口与匹配网络，温度 & MFEM 电磁示例~\cite{mfemem}；综述~\cite{rf}\\
神经响应 & 被动电缆、激活函数、节点模型、强度--时程 & 主动膜与有髓纤维模型、纤维群体、实验阈值 & \cite{rattay,davids}\\
磁滞、热与网格 & 回线面积即损耗；SAR 不是温度；网格必须表达几何 & 磁滞模型、生物传热、网格生成算法 & 专门的课程与文献\\
\bottomrule
\end{tabularx}
\end{table}

\subsection{最后一个问题}
回到全书的第一个问题：头皮上方线圈的磁场在组织里微不足道，组织里为什么仍有足以刺激神经的电场？现在我们能把答案说完整：电场来自磁通的变化，它的形状由组织的边界与电导率决定，界面上的电荷改写了它；要算出它，需要定律、材料、源、边界与观测量这几层信息；要相信算出的数，还需要验证它解对了方程，确认它选对了方程，并说清剩下的不确定性有多大。面对一个新的线圈、一块新的磁轭或一种新的刺激方式，这套做法同样适用——不需要为每一个应用重新发明电磁学。

\begin{takeaways}
\begin{itemize}
\item 从要做的决定出发选观测量，再选近似、表述与检查；求解器、网格与源的表示服从问题。同一台机器上的不同问题，可以是不同的电磁问题。
\item 数值验证问“方程解对了吗”（代码验证、解的验证），物理确认问“方程选对了吗”（与测量比较同一个观测量）。建模假设只有确认能检查；独立表述与模型之间的对照介于两者之间。
\item 每个应用都有必须通过的检查。全书反复出现的失误——峰值与有效值、$10^6$ 与匝数、规范、$|\vb E|$、冻结割线、空气域、对称面、SAR 相加——都不会让程序报错。
\item 误差预算属于“一个问题加一个观测量”：一次只改一个数值环节，独立扰动不确定的输入。算例中 TES 由材料主导（$+22.5\%$、$-15.5\%$），TMS 由分割主导（$\pm6.1\%$），离散误差在 $n=128$ 时都小于 $1\%$。两张网格一致不是误差的上界。
\item 可复现的结果连同模型约定、来源、波形、坐标、数值设置、检查记录与后处理脚本一起保存；单位与归一化跟着数组走，小信号场图带着它的偏置状态。
\item 教学模型到哪里为止，也是结论的一部分。
\end{itemize}
\end{takeaways}

\begin{selfcheck}
\begin{enumerate}
\item 一个场解在网格加密一倍后，所关心的观测量只变化 $0.1\%$，代数残差为 $10^{-13}$。这说明了什么？不说明什么？
\item 下面各属于数值验证、物理确认，还是两者之间？(a) 用制造解测收敛阶；(b) 把气隙磁场与霍尔探头的测量比较；(c) 用能量法和应力张量法分别算同一个力；(d) 把磁路网络与二维场解比较。
\item 在\exref{ex:practice-budget} 的 TES 中，把网格从 $256$ 再加密到 $512$，值得吗？若要缩小 $y$ 的不确定性，该把力气花在哪里？
\item \exref{ex:practice-twomesh} 中，1 个与 2 个单元的结果完全一致。为什么这不能说明网格已经收敛？应当怎样改进网格？
\item 你关心的量是对称面上本应为零的一个磁场分量。用 $|y_{h/2}-y_h|/|y_{h/2}|$ 判断收敛，会出什么问题？
\item 你为偏置磁轭附近的一个匀场线圈导出了一张每安培的小信号场图。除了单位与坐标系，还必须随它保存什么？为什么？
\item 一份报告写着“峰值 SAR 为 $2.5\,\mathrm{W/kg}$”。要使用这个数，你至少还要追问哪几件事？
\end{enumerate}
\end{selfcheck}

\begin{answers}
\item 它说明对\emph{这个}观测量，离散误差与代数误差都很小——这是一项通过了的解的验证。它不说明方程选对了（几何、材料、近似是否正确），不说明其他观测量（例如点峰值）也收敛了，也不说明网格正确地表达了几何。
\item (a) 数值验证（代码验证）；(b) 物理确认；(c) 独立表述，介于两者之间、性质上偏向验证：能抓住实现与推导的错误，但两种方法共用同一组建模假设，对确认没有帮助；(d) 模型之间的对照：用一个更完整、已经验证过的模型检查网络的压缩假设，不是对真实装置的确认。
\item 不值得。$128\to256$ 已经只变 $0.80\%$，而夹杂电导率 $\pm20\%$ 带来 $+22.5\%$、$-15.5\%$ 的变化。力气应花在电导率的数据（测量或有来源的范围），其次是电极模型（最多 $8.5\%$）。
\item 两张网格都把低电导层“丢”了（所有单元中心都在第 1 层），所以它们犯了同一个错误、得到同一个错误答案。应当让网格贴合界面（界面落在节点上），并在加密时检查几何的表示是否随之改进。
\item 分母 $|y_{h/2}|$ 接近零，指标会变得任意大（或在恰为零时没有定义），不能反映真实误差。应改用一个明确的绝对参考尺度 $y_{\rm ref}$，例如该区域中磁场的典型大小 \eqref{eq:practice-indicator}。
\item 非线性偏置状态（磁体剩磁、线圈电流、所用的 B--H 数据及其版本）和对应的切线。小信号场图只是偏置附近的导数 \eqref{eq:biaslinearization}，离开了偏置的说明，就不知道它在哪里有效，也无法复现。
\item 峰值约定还是有效值约定；在多大的质量或区域上平均、怎样处理空气；是否做了时间（脉冲串）平均、占空比多少；源怎样归一化、对应的激发目标是什么；所用的组织参数与模型。约定一错就差 2 倍：同一个“$100\,\mathrm{V/m}$”按峰值读给出 $2.5\,\mathrm{W/kg}$，按有效值读给出 $5\,\mathrm{W/kg}$（\secref{sec:framework-sar}）。
\end{answers}
